\documentclass[12pt]{article}
\usepackage[margin=1in]{geometry}
\usepackage[utf8]{inputenc}
\usepackage{latexsym, amsfonts, enumitem, graphicx, environ,enumitem, fancyhdr, color, amssymb, amsthm, amsmath, mathtools, tikz, nicematrix}
\usepackage{physics}
\usetikzlibrary{patterns}
\usetikzlibrary{decorations.pathreplacing}
\pagestyle{fancy}
\setlength{\headheight}{65pt}
\newenvironment{problem}[2][Problem]{\begin{trivlist}
    \item[\hskip \labelsep {\bfseries #1}\hskip \labelsep {\bfseries #2.}]}{\end{trivlist}}
\DeclareMathOperator{\Ima}{Im}

\usepackage{hyperref}
\usepackage[capitalize]{cleveref}
\usepackage{thmtools}

\newtheorem{thm}{Theorem}
\newtheorem{lem}[thm]{Lemma}
\newtheorem{cl}[thm]{Claim}
\newtheorem{prop}[thm]{Proposition}
\newtheorem{cor}[thm]{Corollary}
\newtheorem{conj}[thm]{Conjecture}
\newtheorem{obstruction}[thm]{Obstruction}

\usepackage{mdframed}

\theoremstyle{definition}
\newtheorem{definition}[thm]{Definition}
\newtheorem{remark}[thm]{Remark}
\newtheorem{question}{Question}
\newtheorem{example}[thm]{Example}
\newtheorem{exercise}[thm]{Exercise}
\newtheorem{properties}[thm]{Properties}

\usepackage{tcolorbox}
\tcbuselibrary{theorems,skins,breakable}

\tcolorboxenvironment{example}{
enhanced jigsaw,colframe=cyan,interior hidden,
breakable,%before skip=10pt,after skip=10pt
}

\tcolorboxenvironment{thm}{
enhanced jigsaw,colframe=red,interior hidden,
breakable,%before skip=10pt,after skip=10pt
}

\tcolorboxenvironment{prop}{
enhanced jigsaw,colframe=red,interior hidden,
breakable,%before skip=10pt,after skip=10pt
}

\tcolorboxenvironment{properties}{
enhanced jigsaw,colframe=red,interior hidden,
breakable,%before skip=10pt,after skip=10pt
}

\tcolorboxenvironment{lem}{
enhanced jigsaw,colframe=red,interior hidden,
breakable,%before skip=10pt,after skip=10pt
}
\tcolorboxenvironment{cor}{
enhanced jigsaw,colframe=red,interior hidden,
breakable,%before skip=10pt,after skip=10pt
}
\tcolorboxenvironment{obstruction}{
enhanced jigsaw,colframe=red,interior hidden,
breakable,%before skip=10pt,after skip=10pt
}

\tcolorboxenvironment{proof}{
enhanced jigsaw, frame hidden,%interior hidden,
breakable,%before skip=10pt,after skip=10pt
}

\renewcommand\qedsymbol{\( \blacksquare \)}

\lhead{Avi Chetlin \\ Assignment 02 \\ 08 September 2026}
\rhead{Dr. Stanislaw Szarek \\ MATH 423 Measure Theory \\ Fall 2026}

\begin{document}
All problems from \textit{Folland, Real Analysis} (2e).

\begin{problem}{1.2.2(e)}
	Show that \( \mathcal{B}_{\mathbb{R}} \) is generated by \( \mathcal{E}_3 = \left\{ (a,b] \mid a < b \right\} \) or \( \mathcal{E}_4 = \left\{ [a,b) \mid a < b \right\} \).
\end{problem}
\begin{proof}
	The proof sketch on p. 22 established that \( \mathcal{E}_j \subset \mathcal{B}_{\mathbb{R}} \) for all \( j \), as well as the inclusion \( \mathcal{B}_{\mathbb{R}} \subset \mathcal{M}(\mathcal{E}_{1}) \).
	Hence, by Lemma 1.1, it suffices to show that \( \mathcal{E}_1 \subset \mathcal{E}_{3,4} \), since then \( \mathcal{B}_{\mathbb{R}} \subset \mathcal{M}(\mathcal{E}_1) \subset \mathcal{M}(\mathcal{E}_{3,4}) \).

	Let \( (a,b) \in \mathcal{E}_1 \) be an arbitrary open interval.
	We can then write
	\begin{equation}
		\label{eq:32}
		(a,b) = \bigcup_{n=1}^{\infty} (a, b - \frac{1}{n}] \in \mathcal{E}_3 = \bigcup_{n=1}^{\infty} [a + \frac{1}{n}, b) \in \mathcal{E}_{4},
	\end{equation}
	and we are done.
\end{proof}

\begin{problem}{1.2.3}
	Let \( \mathcal{M} \) be an infinite \( \sigma \)-algebra. Prove each of the following.
	\begin{enumerate}[label=(\alph*)]
		\item \( \mathcal{M} \) contains an infinite sequence of disjoint sets.

            We proceed in the manner suggested by the instruction document, first proving two lemmas.
          \begin{lem}
            If \( \mathcal{M} \subset \mathcal{P}(X) \) is a \( \sigma \)-algebra and \( A \in \mathcal{M} \), denote
            \begin{equation}
            \label{eq:37}
            \mathcal{M}_A \coloneq \{ E \in \mathcal{M} \mid E \subset A \} = \{ F \cap A \mid F \in \mathcal{M} \}.
            \end{equation}
            Then \( \mathcal{M}_A \) is a \( \sigma \)-algebra on \( A \).
          \end{lem}
		      \begin{proof}
            First, we establish the two containments necessary to show inequality.
            If \( E \in \{ E \in \mathcal{M} \mid E \subset A\} \), then we can clearly write \( E = E \cap A \) by basic set theoretic properties, which immediately shows that \( E \in \{ F \cap A \mid F \in \mathcal{M} \} \).
            Conversely, if \( F \cap A \in \{F \cap A \mid F \in \mathcal{M}\} \), then \( F \cap A \) is a finite union of elements in \( \mathcal{M} \), hence \( F \cap A \in \mathcal{M} \);
            moreover, \( F \cap A \subset A \), showing that the equality holds.

            Now, it is obvious that \( A = A \cap A \in \mathcal{M}_A \) and \( \emptyset = \emptyset \cap A \in \mathcal{M}_A \).
            Let \( \{ B_j \}_1^{\infty} \subset \mathcal{M}_A\) be an arbitrary family, and consider \( B = \bigcup_1^{\infty} B_j \).
            We can write
            \begin{equation}
            \label{eq:38}
            B = \bigcup_{1}^{\infty} (F_j \cap A) = A \cap \bigcup_1^{\infty} F_j,
            \end{equation}
            where \( \{F_j\} \subset \mathcal{M} \).
            Since \( \mathcal{M} \) is a \(\sigma\)-algebra, \( B \in \mathcal{M} \), and since \( B \subset A \), in particular \( B \in \mathcal{M}_A \), so \( \mathcal{M}_A \) is closed under countable unions.
            An identical argument shows that \( \mathcal{M}_A \) is also closed under countable intersections.

            For complements, note than an arbitrary \( G \in \mathcal{M}_A \) can be written as \( G = H \cap A \) for \( H \in \mathcal{M} \).
            Then
            \begin{equation}
            \label{eq:39}
            A \setminus G = A \setminus (H \cap A) = (A \setminus H) \cup (A \setminus A) = (A \setminus H).
          \end{equation}
          Now, \( A \setminus H = A \cap H^{\complement} \) where by \( H^{\complement} \) we mean \( X \setminus H \).
          Furthermore, \( H^{\complement} \in \mathcal{M} \), showing that \( A \setminus G \in \mathcal{M}_A \), and our claim follows.
		      \end{proof}
          \begin{lem}
            \label{lem:2}
            Let \( A, B \in \mathcal{M} \) with \( A \cap B = \emptyset \).
            If both \( \mathcal{M}_A \) and \( \mathcal{M}_B \) are finite, then \( \mathcal{M}_{A \cup B} \) is finite.
            More precisely, \( \# \mathcal{M}_{A \cup B} = \# \mathcal{M}_A \cdot \# \mathcal{M}_B \).
            Consequently, if \( \mathcal{M}_{A \cup B} \) is infinite, at least one of \( \mathcal{M}_A \) or \( \mathcal{M}_B \) is infinite.
          \end{lem}
          \begin{proof}
            Suppose \( A, B \in \mathcal{M} \) are disjoint.
            \( \mathcal{M}_{A \cup B} = \{ E \in \mathcal{M} \mid E \subset A \cup B \} = \{ F \cap (A \cup B) \mid F \in \mathcal{M} \}\).
            Note that we can rewrite this as \( (F \cap A) \cup (F \cap B) \), so that every set \( G \in \mathcal{M}_{A \cup B} \) is the union of corresponding sets \( G_A \in \mathcal{M}_A, G_B \in \mathcal{M}_B \).
            We thus claim that
            \begin{equation}
            \text{card}(\mathcal{M}_{A \cup B}) = \text{card}(\left( \mathcal{M}_A \times \mathcal{M}_B  \right)).
            \end{equation}

            Define \( f: \mathcal{M}_{A \cup B} \rightarrow \mathcal{M}_A \times \mathcal{M}_B \) such that
            \begin{equation}
            \label{eq:40}
            G \mapsto (G_A, G_B) = (F \cap A, F \cap B),
            \end{equation}
            as above.
            Suppose \( f(G_1) = f(G_2) \), so
            \begin{equation}
            \label{eq:41}
            (F_1 \cap A, F_1 \cap B) = (F_2 \cap A, F_2 \cap B).
            \end{equation}
            This obviously implies that
            \begin{align}
            \label{eq:42}
              G_1 &= G_{1,A} \cup G_{1,B} = (F_1 \cap A) \cup (F_1 \cap B)\\
              &= (F_2 \cap A) \cup (F_2 \cap B) = G_{2, A} \cup G_{2, B} = G_2,
            \end{align}
            so \( f \) is injective.

            Let \( (F \cap A, F \cap B) \in \mathcal{M}_A \times \mathcal{M}_B \) be arbitrary.
            Then \( G = (F \cap A) \cup (F \cap B) \) is in \( \mathcal{M}_{A \cup B} \), and \( f(G) = (F \cap A, F \cap B) \), so \( f \) is a bijection.

            Now, for any sets \( S_1, S_2 \), \( \text{card}(\left( S_1 \times S_2 \right)) = \text{card}(S_1 \cdot \#S_2) \), and our claim now follows (since for finite sets, \( \text{card}(S) \coloneq \# S \)).
          \end{proof}
          We may now prove the initial claim.
          \begin{proof}[Proof of 1.2.3(a)]
            We inductively construct such a sequence of disjoint sets, \( \{A_j\}_1^{\infty} \).
            First of all, define \( B_n \) for \( n=0 \) as \( B_0  = X \), and let \( A_1 = \emptyset \).
            Then, have \( B_1 = B_0 \setminus A_1 = X \).
            The associated sub \( \sigma \)-algebra is \( \mathcal{M}_{B_1} = \mathcal{M}_X = \mathcal{M} \) which is, by assumption, infinite.
            In particular, \( \mathcal{M} \neq \{\emptyset, B_1\} \), so there exists some \( E_1 \in \mathcal{M} \) which satisfies
            \begin{equation}
            \label{eq:43}
            \emptyset \subsetneq E_{1} \subsetneq B_1.
            \end{equation}
            Additionally, \( B_1 = (B_1 \setminus E_1) \cup E_1 \).
            By \cref{lem:2}, since \( \mathcal{M} = \mathcal{M}_{B_1} = \mathcal{M}_{(B_1 \setminus E_{1}) \cup E_1} \) is infinite, at least one of \( \mathcal{M}_{B_1 \setminus E_1} \) or \( \mathcal{M}_{E_1} \) are infinite.

            Now, induct on \( n \): Assume we have \( \{ A_j \}_1^n \) which is pairwise disjoint and such that
            \begin{equation}
            \label{eq:45}
            B_n = X \setminus \bigcup_{1}^n A_j
            \end{equation}
            forms an infinite subalgebra \( \mathcal{M}_{B_n} \).
            Then there is some non-empty \( E_{n} \) as in the inductive hypothesis so that
            \begin{equation}
            \label{eq:46}
            \emptyset \subsetneq E_n \subsetneq E_{n-1} \subsetneq \cdots \subsetneq E_1 \subsetneq B_1 = X,
            \end{equation}
            and either \( \mathcal{M}_{B_n \setminus E_{n}} \) or \( \mathcal{M}_{E_{n}} \) is infinite.
            In the former case, put \( A_{n+1} =  E_{n} \), whence \( B_{n+1} = B_n \setminus A_{n+1} \) has an infinite associated subalgebra.
            In the latter case, put \( A_{n+1} = B_n \setminus E_{n} \), and thus \( B_{n+1} = B_n \setminus A_{n+1} = B_n \setminus (B_n \setminus E_{n}) = E_{n} \) again as an infinite associated sub-algebra.

            Since in each case we have guaranteed both \( E_{n} \) and \( B_n \setminus E_{n} \) are non-empty, so too is \( A_{n+1} \).
            Additionally, \( A_{n+1} \subset B_n = X \setminus \bigcup_1^n A_{j} \), so each additional \( A_{n+1} \) is disjoint from the previous \( A_1, \dots, A_n \).

            We have thus constructed an infinite sequence of non--empty, disjoint sets \( \{ A_j \}_1^{\infty} \), as desired.
          \end{proof}
		\item \( \text{card}(\mathcal{M}) \geq \mathfrak{c} \).
		      \begin{proof}
			      By part (a), \( \mathcal{M} \) contains an infinite sequence of disjoint, nonempty sets; call this sequence \( \{ A_n \}_{n \in \mathbb{N}} \).
            Now, define \( f: \mathcal{P}(\mathbb{N}) \rightarrow \mathcal{M} \) by \( N \subset \mathbb{N} \mapsto \{ A_n \}_{n \in N} \).
            Note that since \( \{ A_{n} \} \) is infinite and each set is disjoint, every subsequence \( \{ A_n \}_{n \in N} \) is unique (for unique \( N \subset \mathbb{N} \)).
            That is, \( f \) is an injection, so by definition, \( \text{card}(\mathcal{P}(\mathbb{N})) \leq \text{card}(\mathcal{M}) \).
            By Proposition 0.12 in Folland, \( \text{card}(\mathcal{P}(\mathbb{N})) = \mathfrak{c} \), therefore \( \mathfrak{c} \leq \text{card}(\mathcal{M}) \); finally, by Proposition 0.6,
            \begin{equation}
            \label{eq:33}
            \text{card}(\mathcal{M}) \geq \mathfrak{c}.
            \end{equation}
		      \end{proof}
	\end{enumerate}
\end{problem}

\begin{problem}{1.3.8}
	Show that if \( (X, \mathcal{M}, \mu) \) is a measure space and \( \{E_j\}_1^{\infty} \subset \mathcal{M} \), then\\ \(\mu(\liminf E_j) \leq \liminf \mu(E_j)\).
	Also, \(\mu(\limsup E_j) \geq \limsup \mu(E_j)\), provided that\\ \(\mu(\cup_1^{\infty}) < \infty\).
\end{problem}
\begin{proof}
	To begin with, we recall the definitions of limit superior and inferior for sets and for sequences.
	\begin{alignat}{2}
		\label{eq:1}
		 & \liminf E_j = \bigcup_{k=1}^{\infty} \bigcap_{j=k}^{\infty} E_j               & \qquad
		 & \limsup E_j = \bigcap_{k=1}^{\infty} \bigcup_{j=k}^{\infty} E_j                                                                                                 \\
		 & \liminf x_n = \lim_{k \rightarrow \infty} \inf_{n \geq k}\left\{ x_n \right\}
		 &                                                                               & \limsup x_n = \lim_{k \rightarrow \infty} \sup_{n \geq k} \left\{ x_n \right\}.
	\end{alignat}

	Now, define a sequence of sets \( F_k \coloneq \bigcap_{n=k}^{\infty} E_n \), and note that we have \( F_1 \subset F_2 \subset \dots \subset F_k \), for all \( k \).
	\(\mu\) is a measure, so by continuity from below,
	\begin{equation}
		\label{eq:2}
		\mu(\bigcup_{k=1}^{\infty} F_k) = \lim_{k \rightarrow \infty} \mu(F_k) = \lim_{k\rightarrow \infty} \mu(\bigcap_{n=k}^{\infty} E_n).
	\end{equation}

	Next, since \( \bigcap_{n=k}^{\infty} E_n \subset E_j \) for all \( j \geq k \), we have that \( \mu(\bigcap_{n=k}^{\infty} E_n) \leq \mu(E_j) \).
	In particular, for any \( E_j \subset \{E_{n}\}_k^{\infty} \) such that \(\mu(E_j) = \inf \mu(E_n) + \varepsilon\),
	\begin{equation}
		\label{eq:3}
		\mu(\bigcap_{n=k}^{\infty} E_n) \leq \inf_{j \geq k} \mu(E_n) + \varepsilon.
	\end{equation}
	This holds for all \(\varepsilon > 0\), showing that
	\begin{equation}
		\label{eq:4}
		\mu(\bigcap_{n=k}^{\infty} E_n) \leq \inf_{j \geq k} \mu(E_n).
	\end{equation}
	Taking the limit as \( k \rightarrow \infty \) on both sides shows that
	\begin{equation}
		\label{eq:5}
		\lim_{k\rightarrow \infty} \mu(\bigcap_{n=k}^{\infty} E_n) = \mu(\liminf E_n) \leq \liminf \mu(E_n).
	\end{equation}

	The second claim follows a very similar line of argument.
	Define \( F_k = \bigcup_{n=k}^{\infty} E_{n}\) analogously, and note that \( F_1 \supset F_2 \supset \dots \supset F_k \).
	We are assuming that \( \bigcup_1^{\infty} E_j \) is \( \mu \)-finite, hence, by continuity from above,
	\begin{equation}
		\label{eq:6}
		\mu(\limsup E_{j}) = \mu(\bigcap_{k=1}^{\infty} F_k) = \lim_{k\rightarrow\infty} \mu(F_k).
	\end{equation}
	Moreover, we have for all \( j \geq k \), \( E_j \subset \bigcup_{n=k}^{\infty} E_n \), so just as before,
	\begin{equation}
		\label{eq:7}
		\sup_{j \geq k} \mu(E_j) \leq \mu(\bigcup_{n=k}^{\infty} E_n) = \mu(F_k).
	\end{equation}
	Again taking the limit on both sides, we are left with
	\begin{equation}
		\label{eq:8}
		\mu(\limsup E_j) = \lim_{k\rightarrow \infty} \mu(F_k) \geq \lim_{k\rightarrow \infty} \sup_{j \geq k} \mu(E_k) = \limsup \mu(E_j),
	\end{equation}
	as desired.
\end{proof}

\begin{problem}{1.3.10}
	Given a measure space \( (X, \mathcal{M}, \mu) \) and \( E \in \mathcal{M} \), define \(\mu_E(A) = \mu(A \cap E)\) for \( A \in \mathcal{M} \).
	Show \( \mu_E \) is a measure.
\end{problem}
\begin{proof}
	First of all, since \( A \in \mathcal{M} \), clearly the domain of \(\mu_E\) is \( \mathcal{M} \), which is a \(\sigma\)-algebra.
	Since \( \mathcal{M} \) is a \(\sigma\)-algebra, \( \emptyset \in \mathcal{M} \).
	Let \( E \in \mathcal{M} \) be arbitrary, and consider \(\mu_E(\emptyset) = \mu(\emptyset \cap E) = \mu(\emptyset)\), and since \( \mu \) is a measure, \(\mu_E(\emptyset) = 0\).
	Finally, let \( \{A_j\}_1^{\infty} \) be an arbitrary sequence of disjoint sets in \( \mathcal{M} \), and again let \( E \in \mathcal{M} \) be arbitrary.
	\begin{align}
		\mu_E(\bigcup_1^{\infty} A_j) & = \mu\left(E \cap (\bigcup_{1}^{\infty} A_j) \right) = \mu \left( \bigcup_1^{\infty} (E \cap A_{j}) \right) \\
		\label{eq:disjoint}           & = \sum\limits_1^{\infty} \mu(E \cap A_j) = \sum\limits_1^{\infty} \mu_E(A_j),
	\end{align}
	where \cref{eq:disjoint} follows since if \( \{A_j\} \) are pairwise disjoint, and \( E \cap A_j \subset A_j \), then clearly \( \{ E \cap A_{j }\} \) are disjoint.

\end{proof}

For the following problem, recall that for two sets \( E, M \), we define the symmetric difference as
\begin{equation}
	\label{eq:9}
	E \Delta F \coloneq (E \setminus F) \cup (F \setminus E).
\end{equation}

\begin{problem}{1.3.12}
	Let \( (X, \mathcal{M}, \mu) \) be a finite measure space. Show that
	\begin{enumerate}[label=(\alph*)]
		\item If \( E, F \in \mathcal{M} \) and \(\mu(E \Delta F) = 0\), then \(\mu(E) = \mu(F)\).
		      \begin{proof}
			      Let \( E, F \in \mathcal{M} \) be such that \(\mu(E \Delta F) = 0\).
			      That is, \(\mu \left( (E \setminus F) \cup (F \setminus E) \right) = 0\).
			      Since clearly \( (E \setminus F) \cap (F \setminus E) = \emptyset \), by additivity of \(\mu\) we have
			      \begin{equation}
				      \label{eq:10}
				      \mu( E \Delta F ) = \mu(E \setminus F) + \mu(F \setminus E) = 0,
			      \end{equation}
			      hence \( \mu(E \setminus F) = -\mu(F \setminus E) \).
			      By the non-negativity of \(\mu\), we therefore have that both sides are equal.
			      Moreover, since \( E = (E \setminus F) \cup (E \cap F) \) and \( F = (F \setminus E) \cup (F \cap E) \), we can write
			      \begin{align}
				      \label{eq:14}
				      \mu(E) & = \mu(E \setminus F) + \mu(E \cap F)                 \\
				      \mu(F) & = \mu(F \setminus E) + \mu(F \cap E)= \mu(E \cap F).
			      \end{align}
			      Appealing to the symmetry of intersections, we see that in fact
			      \begin{equation}
				      \label{eq:17}
				      \mu(E) = \mu(F).
			      \end{equation}
		      \end{proof}
		\item Say that \( E \sim F \) if \(\mu(E \Delta F) = 0\); then \( \sim \) is an equivalence relation on \( \mathcal{M} \).
		      \begin{proof}
			      Reflexitivity is immediate, since for any set \( E \), \(\mu(E \Delta E) = \mu(E \setminus E \cup E \setminus E) = \mu(\emptyset \cup \emptyset) = \mu(\emptyset) = 0\).

			      For symmetry, suppose \( E \sim F \) so that \( \mu(E \Delta F) = 0 \). Note that unions are symmetric, and thus in general we have
			      \begin{equation}
				      E \Delta F = (E \setminus F) \cup (F \setminus E) = (F \setminus E) \cup (E \setminus F) = F \Delta E,
			      \end{equation}
			      so \( \mu(E \Delta F) = \mu(F \Delta E) = 0 \) showing \( F \sim E \).

			      Finally, for transitivity, suppose \( E \sim F \) and \( F \sim G \).
			      First we establish that
			      \begin{equation}
				      \label{eq:21}
				      E \Delta G \subset (E \Delta F) \cup (F \Delta G).
			      \end{equation}
			      Suppose that \( x \in E \Delta G \), so \( x \in E \setminus G \) or \( x \in G \setminus E \).
			      If \( x \in E \setminus G \), either \( x \in F \) in which case \( x \in F \setminus G \), or \( x \not\in F \), whence \( x \in E \setminus F \).
			      In either case, \( x \in (E \Delta F) \cup (F \Delta G) \).

			      A symmetric argument holds for when \( x \in G \setminus E \), showing \( E \Delta G \subset (E \Delta F) \cup (F \Delta G) \).

			      By monotonicity and non-negativity,
			      \begin{equation}
				      \label{eq:22}
				      0 \leq \mu(E \Delta G) \leq \mu(E \Delta F) + \mu(F \Delta G) = 0,
			      \end{equation}
			      so \( \mu(E \Delta G) = 0 \), which is to say that \( E \sim G \), as desired.
		      \end{proof}
		\item For \( E, F \in \mathcal{M} \), define \(\rho(E,F) = \mu(E \Delta F)\).
		      Then \( \rho(E,G) \leq \rho(E,F) + \rho(F,G) \), hence \(\rho\) is a metric on \( \mathcal{M} / \sim \).
		      \begin{proof}
			      Let \( \rho \) be defined as above, and let the sets \( E, F, G \in \mathcal{M} \) be arbitrary.

			      In (b) we showed that \( E \Delta G \subset (E \Delta F) \cup (F \Delta G) \).
			      Applying \(\mu\) on both sides, monotonicity of a measure immediately implies
			      \begin{equation}
				      \label{eq:28}
				      \mu(E \Delta G) = \rho(E, G) \leq \mu(E \Delta F) + \mu(F \Delta G) = \rho(E, F) + \rho(F, G).
			      \end{equation}

			      To conclude \(\rho\) is a measure, we show that it is well-defined, positive definite, and symmetric.
			      Symmetry was already shown in (b), and positive definiteness follows directly from the non-negativity of the measure and that \(\mu(E \Delta F) = 0 \iff E \sim F\).

			      For well-definedness, suppose \( [E] = [E^{\prime}], [F] = [F^{\prime}] \), and consider \( \rho(E,F) \) and \( \rho(E^{\prime}, F^{\prime}) \).

            By the triangle inequality we've already proven,
            \begin{equation}
            \label{eq:34}
            \rho(E,F) \leq \rho(E, E^{\prime}) + \rho(E^{\prime}, F) \leq \rho(E, E^{\prime}) + \rho(E^{\prime}, F^{\prime}) + \rho(F^{\prime}, F),
            \end{equation}
            and since \( E \sim E^{\prime}, F \sim F^{\prime} \), we have \(\rho(E, E^{\prime}) = \rho(F, F^{\prime}) = 0\).
            By symmetry, \cref{eq:34} then shows that
            \begin{equation}
            \label{eq:36}
            \rho(E,F) \leq \rho(E^{\prime}, F^{\prime}); \qquad \rho(E^{\prime}, F^{\prime}) \leq \rho(E,F),
            \end{equation}
            hence \(\rho(E, F) = \rho(E^{\prime}, F^{\prime})\).
		      \end{proof}
	\end{enumerate}
\end{problem}
\end{document}
